% ID: SEL-D0054-OB-E01
% ESTADO: cotejada. Segunda pasada visual y borrador compilado revisados.
% Fuente: Selectividad/sel_2011_matematicas_aplicadas/121-EXAMEN 4.pdf
% SHA-256: c70ed065da193c82dd0a02fb4c05069415b0adf53d0f49bba1b0375a19823566
% Página: 2; opción B; ejercicio original 1.

Sea el recinto determinado por las siguientes inecuaciones:
\[
x+y\leq20,\quad3x+5y\leq70,\quad x\geq0,\quad y\geq0.
\]
\begin{enumerate}[label=\alph*)]
\item \textbf{(0.5 puntos)} Razone si el punto de coordenadas $(4.1,11.7)$ pertenece al recinto.
\item \textbf{(1.25 puntos)} Represente dicho recinto y calcule sus vértices.
\item \textbf{(0.75 puntos)} \textquestiondown{}Dónde alcanzará la función $F(x,y)=0.6x+y$ sus valores extremos y cuáles serán éstos?
\end{enumerate}
